% ============================================================================
\section{System-level safeguards}
% ============================================================================

\begin{frame}{Two Layers of Defense}
  \begin{itemize}\itemsep0.5em
    \item[\textcolor{KAred}{$\blacktriangleright$}] \textbf{System level:} what we build around the model, including permissions, sandboxes, approval gates, input handling, and monitors. Most agent safety is won here.
    \item[\textcolor{KAred}{$\blacktriangleright$}] \textbf{Model level:} what we train into the model, including refusal behavior, alignment, and internal monitors.
    \item[\textcolor{KAred}{$\blacktriangleright$}] The working assumption across the field is that the model can be fooled, so the system is designed such that a fooled model still cannot do catastrophic damage.
    \item[\textcolor{KAred}{$\blacktriangleright$}] No single layer is sufficient, so independent layers are stacked. This is defense in depth.
  \end{itemize}
\end{frame}

\begin{frame}{Least Privilege}
  \begin{itemize}\itemsep0.5em
    \item[\textcolor{KAred}{$\blacktriangleright$}] Give the agent the minimum tools, data, and permissions the task needs. This directly counters excessive agency.
    \item[\textcolor{KAred}{$\blacktriangleright$}] \textbf{Tool allowlists:} the agent can call only a fixed, reviewed set of tools, not an open-ended plugin surface.
    \item[\textcolor{KAred}{$\blacktriangleright$}] \textbf{Scoped credentials:} read-only where writing is not needed, one mailbox rather than the whole account, and short-lived tokens over standing access.
    \item[\textcolor{KAred}{$\blacktriangleright$}] With tight scoping, a fully hijacked agent is still bounded by what its credentials allow. The injection may succeed and reach nothing valuable.
  \end{itemize}
\end{frame}

\begin{frame}{Sandboxing}
  \begin{itemize}\itemsep0.5em
    \item[\textcolor{KAred}{$\blacktriangleright$}] Run the agent's actions in an isolated environment so their effects are contained and reversible.
    \item[\textcolor{KAred}{$\blacktriangleright$}] \textbf{Code execution:} a container or VM with no host access, no secrets mounted, and a network that is off by default or allowlisted.
    \item[\textcolor{KAred}{$\blacktriangleright$}] \textbf{Browsing:} an isolated browser profile, separate from the user's real cookies and sessions, so a malicious page cannot ride existing logins.
    \item[\textcolor{KAred}{$\blacktriangleright$}] Sandboxing limits the blast radius. Together with least privilege, it keeps a compromise inside a disposable environment.
  \end{itemize}
\end{frame}

\begin{frame}{Human-in-the-Loop for Consequential Actions}
  \begin{itemize}\itemsep0.5em
    \item[\textcolor{KAred}{$\blacktriangleright$}] Require explicit human approval before any irreversible or high-impact action: payments, deletions, outbound messages, code deployment.
    \item[\textcolor{KAred}{$\blacktriangleright$}] Let low-risk, reversible actions run automatically, and reserve the interrupt for the small set that matters, so approvals stay meaningful.
    \item[\textcolor{KAred}{$\blacktriangleright$}] \textbf{Approval fatigue} is the main failure mode. If the agent asks constantly, users approve reflexively and the gate stops working.
    \item[\textcolor{KAred}{$\blacktriangleright$}] Effective gates show what will happen and why, in plain terms, so the human is actually deciding rather than rubber-stamping.
  \end{itemize}
\end{frame}

\begin{frame}{Handling Untrusted Data, and the Limits of Guardrails}
  \begin{itemize}\itemsep0.5em
    \item[\textcolor{KAred}{$\blacktriangleright$}] \textbf{Mark untrusted input:} wrap retrieved content in delimiters and enforce an instruction hierarchy so tool and web text cannot outrank the user.
    \item[\textcolor{KAred}{$\blacktriangleright$}] \textbf{Guardrail classifiers:} separate models that screen inputs and outputs for injection attempts, policy violations, or data exfiltration.
    \item[\textcolor{KAred}{$\blacktriangleright$}] \textbf{Their limits:} classifiers can themselves be fooled, add latency, and raise false positives that push users to disable them. They lower risk without removing it.
    \item[\textcolor{KAred}{$\blacktriangleright$}] Guardrails work as one layer among several. They sit on top of least privilege and sandboxing, which hold even when the classifier is bypassed.
  \end{itemize}
\end{frame}

\begin{frame}{Separating Control Flow From Untrusted Data}
  \begin{itemize}\itemsep0.45em
    \item[\textcolor{KAred}{$\blacktriangleright$}] Delimiters, privilege ordering, and classifiers can all be talked past. The stronger goal is to make it structurally impossible for untrusted content to choose an action, at the cost of a narrower agent.
    \item[\textcolor{KAred}{$\blacktriangleright$}] \textbf{Dual LLM:} a privileged model holds the tools and never reads untrusted content; a quarantined model reads it and returns values the privileged model only ever handles as opaque variables.
    \item[\textcolor{KAred}{$\blacktriangleright$}] \textbf{Plan-then-execute:} fix the sequence of tool calls from the trusted request alone, before any untrusted data is fetched, so retrieved text can fill in arguments but cannot add steps.
    \item[\textcolor{KAred}{$\blacktriangleright$}] CaMeL implements this with an interpreter that tracks where every value came from and checks a policy before each tool call. It solves 77\% of AgentDojo tasks with provable security, against 84\% undefended.
  \end{itemize}
  \imgsrc{Beurer-Kellner et al., ``Design Patterns for Securing LLM Agents against Prompt Injections,'' \href{https://arxiv.org/abs/2506.08837}{arXiv:2506.08837}; Debenedetti et al., ``Defeating Prompt Injections by Design'' (CaMeL), \href{https://arxiv.org/abs/2503.18813}{arXiv:2503.18813}.}
\end{frame}
